% vim: spl=pt
\begin{exercício}{Propriedade de operadores auto-adjuntos}{ex28}
  Seja \(T \in \bounded(\hilbert)\) um operador auto-adjunto. Mostre que \(0 \leq \norm{T}\unity - T.\)
\end{exercício}
\begin{proof}[Resolução]
  É claro que \(\norm{T} \unity - T\) é auto-adjunto. Pela aplicação espectral, segue que
  \begin{equation*}
    \sigma(\norm{T}\unity - T) = \setc{\norm{T} - \lambda}{\lambda \in \sigma(T)} \subset \setc{\norm{T} - \lambda}{\lambda \in [-\norm{T}, \norm{T}]} \subset [0, 2\norm{T}] \subset \mathbb{R}_+,
  \end{equation*}
  portanto \(\norm{T} \unity - T\) é positivo.
\end{proof}
